\documentclass[11pt,a4paper]{article}
\usepackage[margin=23mm]{geometry}
\usepackage{amsmath,amssymb,amsthm,hyperref}
\newtheorem{theorem}{Theorem}
\newtheorem{lemma}[theorem]{Lemma}
\newcommand{\Del}{\Delta}
\newcommand{\om}{\omega}
\newcommand{\Om}{\Omega}
\title{Divisor concentration and primitive bases}
\author{Human Compression project: working theory for F008}
\date{}
\begin{document}\maketitle
This working note makes explicit the analytic and divisibility arguments
used in the planar Mondrian density theorem of Lebowitz--Lockard and
Tenenbaum~\cite{LLT}. The moment method is that of Hall and
Tenenbaum~\cite{HT}; it is not claimed here as a new method. We give the
prime correlation argument, the passage from squarefree to arbitrary
integers, and the passage from a Dirichlet series to an ordinary mean.
The classical inputs retained below are Chebyshev's bound
$\sum_{p\le z}\log p\ll z$, Mertens' bound
$\sum_{p\le z}p^{-1}\le\log\log z+O(1)$, and elementary Euler summation
for $\zeta$. Constants may depend on the fixed parameters.

\section{Moments of divisors in a short interval}
Write
\[
 \Del(n,u)=\#\{d\mid n:e^u<d\le e^{u+1}\},\quad
 \Del(n)=\sup_u\Del(n,u),\quad
 M_j(n)=\int_{\mathbb R}\Del(n,u)^j\,du.
\]
We use $\log_2x=\log\log x$ and $\log_3x=\log\log\log x$.

\begin{lemma}\label{basic}
For $j\ge1$, $M_1(n)=\tau(n)$,
$M_j(n)\le\tau(n)^j$, and $\Del(n)\le2M_j(n)^{1/j}$.
Also $\Del(ab)\le\tau(b)\Del(a)$ for arbitrary positive integers $a,b$.
For $1<\sigma\le2$,
\[
 \sum_p\frac{\log p}{p^\sigma}<\frac1{\sigma-1}.
\]
\end{lemma}
\begin{proof}
Each divisor contributes an interval of length one to the first integral.
The bound $\Del(n,u)\le\tau(n)$ gives the second assertion.
If a unit interval contains $D$ logarithms of divisors, integration over
the interval of length two covering all their unit windows gives an
integral at least $D$. H\"older's inequality gives
$D\le2^{1-1/j}M_j(n)^{1/j}$.

Every divisor of $ab$ can be assigned uniquely a decomposition $d=vw$
with $v\mid a,w\mid b$: at each prime assign to $v$ the smaller of its
exponent in $d$ and in $a$. For each fixed $w$, the possible $v$ in a
unit logarithmic interval contribute at most $\Del(a)$.

For clarity, the last estimate is the elementary prime-series lemma of
\cite[Lemma 1]{HT}. Put $s=\sigma-1$. Euler summation gives
\[
 \zeta(\sigma)=\frac1s+\frac12+
 \frac{\sigma(\sigma+1)}2
 \int_1^\infty\frac{\{v\}(1-\{v\})}{v^{\sigma+2}}\,dv.
\]
Thus $\zeta(\sigma)>s^{-1}+1/2$. Differentiating the integral, dropping
the negative term and using $\{v\}(1-\{v\})\le1/4$ gives
$-\zeta'(\sigma)<s^{-2}+\sigma/(8(\sigma+1))$.
For $0<s\le1$ the quotient of this bound by $s^{-1}+1/2$ is smaller
than $s^{-1}$. Finally the Euler product gives
$\sum_p(\log p)p^{-\sigma}\le-\zeta'(\sigma)/\zeta(\sigma)$.
\end{proof}

\begin{lemma}[Prime correlation]\label{correlation}
Fix an integer $r\ge1$ and put $a=r+1$. For squarefree $m$, a prime
$p\nmid m$, and an integer $j\ge2$, write
\[
 M_j(m^rp^r)=aM_j(m^r)+W_j(m,p).
\]
Then $W_j(m,p)\ge0$ and
\[
 \sum_{p\nmid m}\frac{\log p}{p}W_j(m,p)
 \ll_r a^j\tau(m^r)M_{j-1}(m^r).
\]
\end{lemma}
\begin{proof}
Set $A_i(u)=\Del(m^r,u-i\log p)$ for $0\le i\le r$.
The identity $\Del(m^rp^r,u)=\sum_i A_i(u)$ shows that the pure
terms in its $j$th power integrate to $aM_j(m^r)$.
Every mixed monomial of degree $j$ is bounded by
$\sum_{i\ne h}A_i^{j-1}A_h$: choose a largest participating factor and
one other participating factor. The sum of all multinomial coefficients
is $a^j$.

For fixed $i\ne h$, change variables to $v=u-i\log p$.
For each divisor $d\mid m^r$, the condition contributing to
$\Del(m^r,v+(i-h)\log p)$ restricts $\log p$ to an interval of
length $1/|i-h|\le1$. Chebyshev's bound and partial summation give a
uniform upper bound for $\sum (\log p)/p$ on every such interval,
including intervals meeting $p=2$. Summing over $d$ therefore bounds
the prime sum of this mixed integral by
$O(\tau(m^r)\int\Del(m^r,v)^{j-1}\,dv)$.
There are $a(a-1)$ ordered pairs, proving the assertion.
\end{proof}

\begin{lemma}[Ordinary means]\label{means}
For fixed $r\ge1,t>0,y>0$, let
$f(n)=\Del(n^r)^t y^{\om(n)}$. Uniformly for $x\ge3$,
\[
 \sum_{n\le x}f(n)\ll_{r,t,y}
 \frac{x}{\log x}\sum_{n\le x}\frac{f(n)}n.
\]
The same assertion holds for $f(n)=y^{\om(n)}$.
\end{lemma}
\begin{proof}
Lemma~\ref{basic} implies
$f(mp^\nu)\le B_\nu f(m)$ with
$B_\nu=\max(1,y)(r\nu+1)^t$, whether or not $p\mid m$.
Since $\log n=\sum_{p^\nu\mid n}\log p$, nonnegativity gives
\[
 \sum_{n\le x}f(n)\log n
 \le\sum_{m\le x}f(m)
       \sum_{p^\nu\le x/m}B_\nu\log p
 \ll x\sum_{m\le x}\frac{f(m)}m.
\]
Indeed the $\nu=1$ sum is $O(x/m)$ by Chebyshev. For $z=x/m$, the
higher powers contribute
\[
 O(\sqrt z(\log z)^{t+1})=O(z),
\]
with the constant covering bounded $z$. For $n>\sqrt x$ we can divide by
$\tfrac12\log x$. For the remaining terms the elementary divisor
bound gives $f(n)\ll n^{1/2}$, hence a total $O(x^{3/4})$, which is
absorbed by $x/\log x$ since the harmonic sum is at least $f(1)=1$.
For the function $y^{\om(n)}$, take $B_\nu=\max(1,y)$ in the same argument.
\end{proof}

\begin{theorem}[Weighted concentration]\label{moment}
Fix an integer $r\ge1$, real $t>0,y>0$, and $\kappa>\sqrt{2t}$.
Put
\[
 \gamma_r(y,t)=y-1+\big((r+1)^t y-y-t\big)^+.
\]
Then, as $x\to\infty$,
\[
 \sum_{n\le x}\Del(n^r)^t y^{\om(n)}
 \ll x(\log x)^{\gamma_r(y,t)}
 \exp\{\kappa\sqrt{\log_2x\log_3x}\}.
\]
\end{theorem}
\begin{proof}
Put $a=r+1$. First suppose $y\ge t/(a^t-1)$, and choose an integer
$j\ge\max(2,t)$. Define the absolutely convergent series
\[
 L(\sigma)=\sum_{m\ {\rm squarefree}}
 \frac{M_j(m^r)^{t/j}y^{\om(m)}}{m^\sigma},\qquad 1<\sigma\le2.
\]
Absolute convergence, also after differentiation, follows from
$M_j(m^r)\le a^{j\om(m)}$ and the corresponding Euler product.
Differentiation singles out one prime. Concavity of $v^{t/j}$ separates
the pure and mixed terms of Lemma~\ref{correlation}. H\"older in the
prime sum, with weights $(\log p)/p^\sigma$, and Lemma~\ref{basic} give
\[
 -L'(\sigma)\le
 \frac{ya^{t/j}L(\sigma)}{\sigma-1}
 +\frac{yC_r^{t/j}a^t}{(\sigma-1)^{1-t/j}}
 \sum_m^*\frac{y^{\om(m)}
 [\tau(m^r)M_{j-1}(m^r)]^{t/j}}{m^\sigma}.
\]
Here and below a star means squarefree. Interpolation between $M_1$
and $M_j$ gives
\[
 \tau(m^r)M_{j-1}(m^r)
 \le a^{j\om(m)/(j-1)}M_j(m^r)^{(j-2)/(j-1)}.
\]
For $j=2$ this is the equality $\tau M_1=\tau^2$.
Another application of H\"older, now to the $m$ sum, therefore yields
\[
 -L'(\sigma)\le\frac{bL(\sigma)}s
 +\frac{BL(\sigma)^{(j-2)/(j-1)}}{s^{1+\alpha/(j-1)}},
 \quad s=\sigma-1,
\]
where
\[
 b=ya^{t/j},\qquad \alpha=a^ty-t+t/j,\qquad
 B=yC_r^{t/j}a^t2^{a^ty/(j-1)}.
\]
We used
$\sum_m^*(a^ty)^{\om(m)}m^{-\sigma}\le\zeta(\sigma)^{a^ty}
\le(2/s)^{a^ty}$. The coefficient $B$ is bounded independently of $j$.

The difference $\alpha-b$ needs a positive lower bound. At the smallest
allowed $y$ it equals
\[
 \frac tj-\frac{t(a^{t/j}-1)}{a^t-1}.
\]
Since $j(a^{t/j}-1)\le2(a^{t/2}-1)$ for $j\ge2$, convexity gives
\[
 \alpha-b\ge\frac tj\frac{a^{t/2}-1}{a^{t/2}+1}=\frac{c_t}{j}>0.
\]
Increasing $y$ only increases this difference. This replaces the overly
strong displayed bound $\alpha-b\ge t/j$ in the source proof.

Set $Y=L^{1/(j-1)}$. The preceding differential inequality is linear
in $Y$. Multiplication by $s^{b/(j-1)}$ and integration from $s$ to one
give
\[
 Y(1+s)\le
 \left(Y(2)+\frac B{\alpha-b}\right)s^{-\alpha/(j-1)}.
\]
Since $L(2)\le\zeta(2)^{a^ty}<2^{a^ty}$, it follows that
\[
 L(1+s)\le(Cj)^j s^{-a^ty+t-t/j}.
\]

We now remove the squarefree restriction. Write $n=uv$, where $u$
is the product of the primes appearing to exponent one in $n$ and $v$
contains the remaining prime powers. Then $u$ is squarefree, $v$ is
powerful, and $(u,v)=1$. Lemma~\ref{basic} gives
$\Del(n^r)\le\tau(v^r)\Del(u^r)$.
The series
\[
 \sum_{v\ {\rm powerful}}\frac{\tau(v^r)^t y^{\om(v)}}{v^\sigma}
 =\prod_p\left(1+y\sum_{\nu\ge2}
           \frac{(r\nu+1)^t}{p^{\nu\sigma}}\right)
\]
is bounded uniformly for $\sigma\ge1$. Thus the full Dirichlet series
of $f(n)=\Del(n^r)^t y^{\om(n)}$ is at most a fixed constant times
$L(\sigma)$, using $\Del(u^r)^t\le2^tM_j(u^r)^{t/j}$.
Lemma~\ref{means}, with $\sigma=1+1/\log x$ and
$n^{-1}\le e n^{-\sigma}$ for $n\le x$, now gives
\[
 \sum_{n\le x} f(n)\ll
 x(\log x)^{a^ty-t-1}
 \exp\{(t/j)\log_2x+j\log(Cj)\}.
\]
For sufficiently large $x$, take
\[
 j=\lfloor\sqrt{2t\log_2x/\log_3x}\rfloor.
\]
The last exponent is then
$(\sqrt{2t}+o(1))\sqrt{\log_2x\log_3x}$.

Finally suppose $0<y<t/(a^t-1)$. The function $T\mapsto T/(a^T-1)$
is continuous and strictly decreasing to zero. Hence there is a $T>t$
with $y=T/(a^T-1)$.
Apply the case just proved with $T$ and then H\"older to
\[
 \sum_{n\le x}\Del(n^r)^t y^{\om(n)}
 \le\left(\sum_{n\le x}\Del(n^r)^T y^{\om(n)}\right)^{t/T}
       \left(\sum_{n\le x}y^{\om(n)}\right)^{1-t/T}.
\]
The final mean is $O(x(\log x)^{y-1})$: its Euler product is at most
$C_y\zeta(\sigma)^y$, and Lemma~\ref{means} applies. At the chosen
threshold both logarithmic powers are $y-1$. The exponential constant
can be chosen arbitrarily close to $t\sqrt{2/T}<\sqrt{2t}$.
This proves the remaining case.
\end{proof}

\section{Exceptional divisors and sets of multiples}
For a set $S$ of positive integers, let ${\cal M}(S)$ be the set of
integers divisible by at least one member of $S$.

\begin{lemma}[Excess prime factors]\label{tails}
Fix $\varepsilon>0$ and let
\[
 E=\{b>1:\Om(b)>(1+\varepsilon)\log_2b\}.
\]
There is $\eta>0$ such that the upper density of
${\cal M}(E\cap[N,\infty))$ is $O((\log N)^{-\eta})$.
\end{lemma}
\begin{proof}
Choose $1<z<2$ sufficiently close to one that
$\eta=(1+\varepsilon)\log z-(z-1)>0$.
Let $n_Y$ denote the part of $n$ supported on primes at most $Y$.
For every $X,Y\ge2$, expansion as a divisor sum gives
\[
 \frac1X\sum_{n\le X}z^{\Om(n_Y)}
 \le\prod_{p\le Y}\frac{1-p^{-1}}{1-zp^{-1}}
 \ll_z (\log Y)^{z-1}.
\]
Indeed the divisor coefficients at $p^\nu$ are
$(z-1)z^{\nu-1}\ge0$; replace each $\lfloor X/d\rfloor/X$ by $1/d$.
Taking logarithms of the product and using Mertens' bound proves the
last inequality; the terms of degree at least two sum absolutely.

If $e^{e^j}\le b<e^{e^{j+1}}$, $b\mid n$ and $b\in E$, then
$\Om(n_{e^{e^{j+1}}})>(1+\varepsilon)j$. Markov's inequality
bounds the proportion of such $n\le X$ by $O(e^{-\eta j})$,
uniformly in $X$ and $j$. Summing over
$j\ge\lfloor\log_2N\rfloor$ proves the assertion, with a constant
covering the first incomplete interval.
\end{proof}

\begin{lemma}[Density below one]\label{multiples}
Suppose $S\subset\mathbb N\setminus\{1\}$ and the upper density of
${\cal M}(S\cap[N,\infty))$ tends to zero. Then ${\cal M}(S)$ has
an asymptotic density strictly smaller than one. In particular this
holds if $\sum_{b\in S}b^{-1}<\infty$.
\end{lemma}
\begin{proof}
The multiples of a finite set have a periodic density. Their difference
from ${\cal M}(S)$ lies in a tail whose upper density tends to zero.
Their densities consequently converge, and the upper and lower densities
of ${\cal M}(S)$ coincide. The same argument applies to every tail of $S$.

We recall the elementary finite form of Behrend's inequality. For finite
sets $A,B$, the densities of their complements satisfy
\[
 d({\cal M}(A\cup B)^c)\ge
 d({\cal M}(A)^c)d({\cal M}(B)^c).
\]
To see this, sample an integer uniformly modulo the least common multiple
of their entries. Its truncated prime valuations are independent. Each
event of avoiding a prescribed collection of divisors is decreasing in
those valuations. Two decreasing functions of one variable have
nonnegative covariance, by expanding
$\tfrac12\mathbb E[(f(X)-f(X'))(g(X)-g(X'))]\ge0$.
Induction on the independent coordinates proves the same assertion for
the product measure, giving the inequality.

Choose $N$ so that the density of ${\cal M}(S\cap[N,\infty))$ is
less than $1/2$. The finite set $A=S\cap[1,N)$ has complement of
positive density: the residue class one modulo its least common multiple
avoids every entry. Apply the finite inequality to $A$ and larger finite
parts of the tail, then pass to their densities. The complement of
${\cal M}(S)$ has density at least $d({\cal M}(A)^c)/2>0$.
Finally the upper density of a tail is at most its reciprocal sum.
\end{proof}

\paragraph{Boundary of this note.}
Theorem~\ref{moment} recovers source Theorem 4.1 when $r=2,t\ge1$, and its
unweighted-power specialization for general $r$. The proof also covers
$0<t<1$: it only uses $0<t/j\le1$, and no step requires $t\ge1$.
This additional range still awaits distinct review. It does not establish
the source's full multiplicative-weight remark or its proposed improved
exponential error. The three-dimensional geometric application remains
separate. This working note has not received distinct review and adds
no completed-original, conjecture or formal-verification credit.

\begin{thebibliography}{9}
\bibitem{LLT} N. Lebowitz--Lockard and G. Tenenbaum,
\emph{On solutions to the Mondrian Puzzle}, retained thirteen-page
preprint dated 28 August 2026.
\bibitem{HT} R. R. Hall and G. Tenenbaum,
\emph{The average orders of Hooley's $\Delta_r$-functions, II},
Compositio Math. \textbf{60} (1986), 163--186.
\end{thebibliography}
\end{document}
